\section{Discussion and Conclusion}
This manuscript places three aspects of sparse state preparation in a common framework: simultaneous resource bounds within the final output width, a composable lifting from small to superlinear support, and a posteriori target-error certification for finite-angle circuits. The theoretical bounds do not treat multicontrolled gates, broadcasts, or repeated reflections as free, and do not hide data-loading costs in an uncharged query model.

Zero additional ancillas does not mean that no intermediate step may use workspace. The requirement is that all workspace comes from the final output register and that its clean-input promises and restoration times are justified explicitly. The width-$d$ feature computation borrows unfinished result locations; reconstruction with sufficient space also permits written outputs to be borrowed as dirty workspace and restored exactly. Householder conjugation turns a fixed unitary that need only prepare one vector correctly into a full-space reflection, avoiding the destruction of earlier amplitudes that can result from simply concatenating batch preparation circuits.

The construction gives upper bounds for all $2\le d\le2^n$, with worst-case matching claimed only in each fixed polynomial-sparsity window. Finite implementation constants may still be large. A finite angle specification is not a Clifford+$T$ gate sequence: the phase and error contracts of single-rotation synthesis must be integrated separately\cite{ross2014}. Hardware connectivity, error-correction overhead, and noise are likewise not determined by the present bounds.

A posteriori certification guarantees every accepted instance, not that the double-precision frontend can satisfy an arbitrarily small requested error. Extremely high precision or mixed dynamic ranges may be explicitly rejected; the target must not be changed, support discarded, or the budget loosened to avoid such rejection. Approximation optimality, elimination fill-in for general sparse isometries, and superiority to external practical synthesizers lie outside the established scope.

The main questions for independent review are the all-size arguments for the divisor-closed evaluation basis and future-location recycling, the full workspace budget for two-batch reconstruction, the interfaces and phases of complex batched reflections, and whether the closest published results imply the complete construction through a direct reduction. Finite experiments cannot replace these proofs or establish historical uniqueness.
